% Unit 111 English translation of chapter8.tex lines 925--1058.
% Source: Methods of Algebra, Volume 2, commit
% 9a5803ff77dd3257484cb177f851a73770a59dd3 (CC BY 4.0).
% stable-unit-id: o014.aljabr2.chapter8.dold-kan-b
% Coverage: the remainder of Section 8.5, from the Dold--Kan theorem through
% the remark on cochain-complex and cosimplicial versions; Unit 112 begins Section 8.6 at line 1059.

\hypertarget{o014.aljabr2.chapter8.dold-kan-b}{ }
% segment-id: o014.aljabr2.chapter8.dold-kan-b.t001
\begin{theorem}[A.\ Dold, D.\ Kan]\label{prop:Dold-Kan}
	\index{Dold--Kan correspondence}
	For every additive category $\mathcal{A}$, the functor
	$\Gamma:\cate{Ch}_{\geq0}(\mathcal{A})\to\cate{s}\mathcal{A}$ is fully
	faithful.
	\begin{itemize}
		% segment-id: o014.aljabr2.chapter8.dold-kan-b.i001
		\item If $\mathcal{A}$ is also a Karoubi category as mentioned in
		\S\ref{sec:direct-sum}, meaning that every idempotent has a kernel,
		then $\Gamma$ is an equivalence of categories.
		% segment-id: o014.aljabr2.chapter8.dold-kan-b.i002
		\item If, in addition, $\mathcal{A}$ is an abelian category (and hence
		also a Karoubi category), then $\mathrm{N}$ is a quasi-inverse functor
		to $\Gamma$, and
		% segment-id: o014.aljabr2.chapter8.dold-kan-b.g001
		$\begin{tikzcd}
			\Gamma: \cate{Ch}_{\geq 0}(\mathcal{A}) \arrow[shift left, r] & \cate{s}\mathcal{A} : \mathrm{N} \arrow[shift left, l]
		\end{tikzcd}$
		is an adjoint equivalence. The unit and counit morphisms of the
		adjunction are described in Lemma~\ref{prop:GammaN-adjunction}.
	\end{itemize}
\end{theorem}
% segment-id: o014.aljabr2.chapter8.dold-kan-b.pf001
\begin{proof}
	Consider the functor
	% segment-id: o014.aljabr2.chapter8.dold-kan-b.q001
	\[ \tilde{h}_{\mathcal{A}}: \mathcal{A} \to \tilde{\mathcal{A}}^\wedge := \cate{Ab}^{(\mathcal{A}^{\opp})}, \quad Y \mapsto \Hom_{\mathcal{A}}(\cdot, Y). \]
	Its composite with the forgetful functor
	$\cate{Ab}\to\cate{Set}$ is the Yoneda embedding
	$h_{\mathcal{A}}:\mathcal{A}\to\mathcal{A}^\wedge$ recalled in
	\S\ref{sec:Yoneda}. The functor $\tilde{h}_{\mathcal{A}}$ is fully
	faithful. This is simply the $\cate{Ab}$-enriched version of the Yoneda
	lemma, but it can also be deduced from the original version. For every
	$Y_1,Y_2\in\Obj(\mathcal{A})$, the composite map
	% segment-id: o014.aljabr2.chapter8.dold-kan-b.q002
	\[ \Hom_{\mathcal{A}}(Y_1, Y_2) \to \Hom_{\tilde{\mathcal{A}}^\wedge}\left(\tilde{h}_{\mathcal{A}}(Y_1), \tilde{h}_{\mathcal{A}}(Y_2) \right) \to \Hom_{\mathcal{A}^\wedge}\left( h_{\mathcal{A}}(Y_1), h_{\mathcal{A}}(Y_2)\right) \]
	is known to be bijective, while the second map is plainly injective.
	Therefore both maps are bijective.

	Proposition~\ref{prop:functor-cat-Abel} shows that
	$\tilde{\mathcal{A}}^\wedge$ is naturally an abelian category. Use the
	same symbol $\tilde{h}_{\mathcal{A}}$ for the induced functors on chain
	complexes and simplicial objects, and consider the diagram
	% segment-id: o014.aljabr2.chapter8.dold-kan-b.e001
	\begin{equation*}\begin{tikzcd}
		\cate{Ch}_{\geq 0}(\mathcal{A}) \arrow[d, "{\tilde{h}_{\mathcal{A}}}"'] \arrow[r, "\Gamma"] & \cate{s}\mathcal{A} \arrow[d, "{\tilde{h}_{\mathcal{A}}}"] \\
		\cate{Ch}_{\geq 0}(\tilde{\mathcal{A}}^\wedge) \arrow[r, "\Gamma"'] & \cate{s}\tilde{\mathcal{A}}^\wedge.
	\end{tikzcd}\end{equation*}
	This diagram commutes up to a canonical isomorphism. The observation in the
	preceding paragraph shows that both vertical arrows are fully faithful,
	while applying Lemma~\ref{prop:DK-Ab} objectwise shows that the second row
	is an equivalence of categories. Hence $\Gamma$ in the first row is fully
	faithful.

	For any functor $F:\mathcal{C}\to\mathcal{C}'$, an object
	$X'\in\Obj(\mathcal{C}')$ is said to lie in the \emph{essential image} of
	$F$ if there is an $X\in\Obj(\mathcal{C})$ such that $X'\simeq FX$.
	Together with the description in Lemma~\ref{prop:DK-Ab} of the
	quasi-inverse to
	$\cate{Ch}_{\geq0}(\cate{Ab})\to\cate{s}\cate{Ab}$, this yields the
	following conclusion: $X\in\Obj(\cate{s}\mathcal{A})$ lies in the
	essential image of $\Gamma$ if and only if
	$(\mathrm{N}\tilde{h}_{\mathcal{A}}(X))_n$ lies in the essential image of
	$\tilde{h}_{\mathcal{A}}:\mathcal{A}\to
	\tilde{\mathcal{A}}^\wedge$ for every $n\geq0$.

	Notice that
	$\tilde{h}_{\mathcal{A}}:\mathcal{A}\to\tilde{\mathcal{A}}^\wedge$
	preserves finite direct sums. If $\mathcal{A}$ is a Karoubi category, the
	essential image of $\mathcal{A}\to\tilde{\mathcal{A}}^\wedge$ is
	therefore closed under taking direct summands. Since
	$(\mathrm{N}\tilde{h}_{\mathcal{A}}(X))_n$ is a direct summand of
	$(\Gamma\mathrm{N}\tilde{h}_{\mathcal{A}}(X))_n\simeq
	\tilde{h}_{\mathcal{A}}(X)_n$,\footnote{Translator's note: the final term
	is printed as $h_{\mathcal{A}}(X)_n$ in the source. The other objects in
	this isomorphism are $\cate{Ab}$-valued presheaves, so the correctly typed
	term is $\tilde{h}_{\mathcal{A}}(X)_n$.} in this case $\Gamma$ is fully
	faithful and essentially surjective, hence an equivalence.

	Finally, the adjoint equivalence theorem
	\cite[Theorem 2.6.12]{Li1} ensures that a quasi-inverse to
	$\Gamma:\cate{Ch}_{\geq0}(\mathcal{A})\to\cate{s}\mathcal{A}$ can always
	be made into its right adjoint. If $\mathcal{A}$ is an abelian category,
	Lemma~\ref{prop:GammaN-adjunction} has already shown that $\mathrm{N}$ is
	the right adjoint of $\Gamma$, and hence is also its quasi-inverse. This
	proves the result.
\end{proof}

% segment-id: o014.aljabr2.chapter8.dold-kan-b.cv001
\begin{convention}
	In view of this result, if $\mathcal{A}$ is a Karoubi category, we may
	choose a quasi-inverse functor to
	$\Gamma:\cate{Ch}_{\geq0}(\mathcal{A})\to\cate{s}\mathcal{A}$ and
	naturally denote it by $\mathrm{N}$.
\end{convention}

% segment-id: o014.aljabr2.chapter8.dold-kan-b.p001
Although $\mathrm{N}X$ was initially defined as a subobject of
$\mathrm{C}X$, it can also be understood as a quotient object; the latter
viewpoint is often more convenient. We now explain this.

% segment-id: o014.aljabr2.chapter8.dold-kan-b.p002
Let $\mathcal{A}$ be a Karoubi category and
$X\in\Obj(\cate{s}\mathcal{A})$. For every $n\geq0$, define the morphism
$\Psi_n$ in $\mathcal{A}$ as the one characterized by the commutativity of
the following diagram for every $0\leq j<n$:
% segment-id: o014.aljabr2.chapter8.dold-kan-b.g002
\[\begin{tikzcd}
	\bigoplus_{0 \leq i < n} X_{n-1} \arrow[r, "\Psi_n"] & X_n \\
	X_j \arrow[hookrightarrow, u, "\text{direct summand}"] \arrow[ru, "{s_j}"'] &
\end{tikzcd}\]
Thus taking $\Coker(\Psi_n)$ (that is, the coequalizer of $\Psi_n$ and $0$)
intuitively removes every degenerate part from $X_n$, at least when
$\mathcal{A}$ is an abelian category. The next result not only guarantees
that $\Coker(\Psi_n)$ exists, but also identifies it with
$(\mathrm{N}X)_n$ via the natural embedding
$u:\mathrm{N}X\to\mathrm{C}X$ in Proposition~\ref{prop:Dold-Kan-u}. This
gives another perspective on $\mathrm{N}X$.

% segment-id: o014.aljabr2.chapter8.dold-kan-b.pr001
\begin{proposition}\label{prop:Dold-Kan-v}
	In the situation above, $\Coker(\Psi_n)$ exists and is canonically
	isomorphic to $(\mathrm{N}X)_n$. Consider the corresponding composite
	$v_n:X_n\twoheadrightarrow\Coker(\Psi_n)\rightiso(\mathrm{N}X)_n$.
	Then $(v_n)_n$ gives a morphism
	$v:\mathrm{C}X\to\mathrm{N}X$ satisfying
	$vu=\identity_{\mathrm{N}X}$ for the natural embedding
	$u:\mathrm{N}X\hookrightarrow\mathrm{C}X$; this construction is
	functorial in $X$.
\end{proposition}
% segment-id: o014.aljabr2.chapter8.dold-kan-b.pf002
\begin{proof}
	By Theorem~\ref{prop:Dold-Kan}, we may assume that $X=\Gamma A$, where
	$A\in\Obj(\cate{Ch}_{\geq0}(\mathcal{A}))$. In this case, $\Psi_n$ can
	be written as
	% segment-id: o014.aljabr2.chapter8.dold-kan-b.q003
	\[ \bigoplus_{0 \leq i < n} \bigoplus_{t: [n-1] \twoheadrightarrow [k]} A_k \to \bigoplus_{t': [n] \twoheadrightarrow [k] } A_k. \]
	From Definition~\ref{def:DK-Gamma} of the functor $\Gamma$, we see that
	$\Psi_n$ acts on the $(i,t)$-summand on the left as follows. Form the
	composite $[n]\xrightarrow{\mathrm{s}^i}[n-1]\xrightarrow{t}[k]$, which
	is still an order-preserving surjection, and denote it by $t'$; then $A_k$
	is mapped identically to the $A_k$-summand on the right corresponding to
	$t'$.

	Observe that $t':[n]\twoheadrightarrow[k]$ arises from such a pair $(i,t)$
	if and only if $k<n$. Using the $\eta$ of
	Lemma~\ref{prop:DK-N-unit}, define $v_n$ to be the composite
	$(\Gamma A)_n\xrightarrow{\text{projection}}A_n
	\xrightarrow[\sim]{\eta}(\mathrm{N}\Gamma A)_n$. The discussion above
	shows that
	% segment-id: o014.aljabr2.chapter8.dold-kan-b.q004
	\[ [\text{quotient}: (\Gamma A)_n \to \Coker \Psi_n] \simeq [v_n: (\Gamma A)_n \twoheadrightarrow (\mathrm{N}\Gamma A)_n ] . \]

	Expanding the definitions immediately gives
	$v_nu_n=\identity_{(\mathrm{N}\Gamma A)_n}$. We next show that $(v_n)_n$
	gives a morphism $v$ in $\cate{Ch}_{\geq0}(\mathcal{A})$. This is
	equivalent to saying that the outer rectangle of the following diagram
	commutes:
	% segment-id: o014.aljabr2.chapter8.dold-kan-b.g003
	\[\begin{tikzcd}[column sep=large]
		(\Gamma A)_n \arrow[r, "{\sum_i (-1)^i d^{\Gamma A}_i}"] \arrow[d, "\text{projection}"'] & (\Gamma A)_{n-1} \arrow[d, "\text{projection}"] \\
		A_n \arrow[r, "{\partial^A_n}"] \arrow[d, "\eta"'] & A_{n-1} \arrow[d, "\eta"] \\
		(\mathrm{N}\Gamma A)_n \arrow[r, "{\partial^{\mathrm{N}\Gamma A}_n}"'] & (\mathrm{N}\Gamma A)_{n-1}.
	\end{tikzcd}\]
	The lower part commutes because $\eta$ is a morphism, while the
	commutativity of the upper part follows by directly applying the definition
	of $d_i^{\Gamma A}$ (take $f=\mathrm{d}^i$ in
	Definition~\ref{def:DK-Gamma}). Finally, the functoriality of $v$ in $X$
	is clear.
\end{proof}

% segment-id: o014.aljabr2.chapter8.dold-kan-b.p003
Our next goal is to compare $\mathrm{C}$ and $\mathrm{N}$ more precisely in
the case of an abelian category. There is a notion of homotopy for morphisms
between chain complexes; see Definition~\ref{def:homotopy} and
Remark~\ref{rem:Hom-chain-cplx}. A morphism of chain complexes is called a
\emph{quasi-isomorphism} if it induces isomorphisms on homology.

% segment-id: o014.aljabr2.chapter8.dold-kan-b.t002
\begin{theorem}\label{prop:Dold-Kan-N-C}
	Let $\mathcal{A}$ be an abelian category and
	$X\in\Obj(\cate{s}\mathcal{A})$. Then
	$u:\mathrm{N}X\to\mathrm{C}X$ from Lemma~\ref{prop:Dold-Kan-u} and
	$v:\mathrm{C}X\to\mathrm{N}X$ from Proposition~\ref{prop:Dold-Kan-v}
	are both quasi-isomorphisms.
\end{theorem}
% segment-id: o014.aljabr2.chapter8.dold-kan-b.pf003
\begin{proof}
	We know that $vu=\identity_{\mathrm{N}X}$, so it is enough to prove that
	$v$ is a quasi-isomorphism. Since $v$ has a right inverse, it is an
	epimorphism. The long exact sequence of chain complexes
	(Proposition~\ref{prop:long-exact-sequence-ses}) reduces the problem to
	proving that $\Ker(v)$ is acyclic. By Theorem~\ref{prop:Dold-Kan}, we may
	henceforth assume that $X=\Gamma A$, where
	$A\in\Obj(\cate{Ch}_{\geq0}(\mathcal{A}))$. Therefore
	% segment-id: o014.aljabr2.chapter8.dold-kan-b.q005
	\begin{gather*}
		(\mathrm{C}X)_n = \bigoplus_{t: [n] \twoheadrightarrow [k]} A_k, \\
		\partial_n : (\mathrm{C}X)_n \to (\mathrm{C}X)_{n-1}.
	\end{gather*}

	For $n\geq0$ and $i\in\Z$, let $(\mathrm{C}^{\leq i}X)_n$ denote the
	subobject of $(\mathrm{C}X)_n$ cut out by the direct summands satisfying
	% segment-id: o014.aljabr2.chapter8.dold-kan-b.q006
	\[ \exists \; j, \quad \max\{n-i, 1\} \leq j \leq n, \quad t(j) = t(j-1). \]
	We have $i\leq i'\implies(\mathrm{C}^{\leq i}X)_n\subset
	(\mathrm{C}^{\leq i'}X)_n$.

	Since $v_n$ is simply the projection onto the summand
	$t=\identity_{[n]}$, if $i\geq n-1$ then
	$(\mathrm{C}^{\leq i}X)_n=\Ker(v_n)$. We now show that
	% segment-id: o014.aljabr2.chapter8.dold-kan-b.e002
	\begin{equation}\label{eqn:Dold-Kan-N-C-aux0}
		\left( (\mathrm{C}^{\leq i} X)_n \right)_{n \geq 0} \;\text{forms a chain subcomplex of}\; \mathrm{C}X \;\text{denoted by}\; \mathrm{C}^{\leq i} X.
	\end{equation}

	To prove this, fix a map $t:[n]\twoheadrightarrow[k]$ and an index
	$\max\{n-i,1\}\leq j\leq n$ satisfying the condition above. From now on
	write $d_i=d_i^X$ and $s_j=s_j^X$. For every $0\leq h\leq n$, take
	$m=n-1$ and $f=\mathrm{d}^h:[n-1]\hookrightarrow[n]$ in the diagram
	\eqref{eqn:Dold-Kan-epi-mono}.
	\begin{itemize}
		% segment-id: o014.aljabr2.chapter8.dold-kan-b.i003
		\item Suppose $h\notin\{j-1,j\}$ and set
		$j':=(\mathrm{d}^h)^{-1}(j)$. Then
		$u:[n-1]\twoheadrightarrow[l]$ satisfies $u(j')=u(j'-1)$, and it is
		easy to see that $n-1-i\leq j'\leq n-1$. In this case,
		$(-1)^hd_h:X_n\to X_{n-1}$ maps $A_k$ to the direct summand of
		$(\mathrm{C}^{\leq i}X)_{n-1}$ determined by $u$.
		% segment-id: o014.aljabr2.chapter8.dold-kan-b.i004
		\item In the situation of the preceding item, if also
		$h\geq n-i-1$, the range of $j'$ can be sharpened to
		$n-i\leq j'\leq n-1$. In other words, the image of $A_k$ then lies
		in $(\mathrm{C}^{\leq i-1}X)_{n-1}$.
		% segment-id: o014.aljabr2.chapter8.dold-kan-b.i005
		\item For $h\in\{j-1,j\}$, verify that
		$t\mathrm{d}^{j-1}=t\mathrm{d}^j$. It follows that
		$d_{j-1},d_j:X_n\rightrightarrows X_{n-1}$ agree on the direct
		summand $A_k$ corresponding to $t$, so $(-1)^{j-1}d_{j-1}$ and
		$(-1)^jd_j$ cancel on that summand.
	\end{itemize}
	This proves \eqref{eqn:Dold-Kan-N-C-aux0}. The second and third items also
	give the congruence
	% segment-id: o014.aljabr2.chapter8.dold-kan-b.e003
	\begin{equation}\label{eqn:Dold-Kan-N-C-aux}
		\sum_{h=0}^n (-1)^h d_h \equiv \sum_{h=0}^{n-i-2} (-1)^h d_h \pmod{ \mathrm{C}^{\leq i-1} X }.
	\end{equation}

	Thus it is enough to prove that $\mathrm{C}^{\leq i}X$ is acyclic. Set
	% segment-id: o014.aljabr2.chapter8.dold-kan-b.q007
	\[ D:=\mathrm{C}^{\leq i}X/\mathrm{C}^{\leq i-1}X. \]
	Notice that $\mathrm{C}^{<-1}X=0$. The long exact sequence of chain
	complexes then reduces the problem to proving inductively that $D$ is
	acyclic, for $i\geq0$.

	To prove that $D$ is acyclic, consider the family of morphisms
	% segment-id: o014.aljabr2.chapter8.dold-kan-b.q008
	\begin{gather*}
		(-1)^{n-i-1}s_{n-i-1}:X_n\to X_{n+1}, \quad n\geq i+1.
	\end{gather*}
	A direct verification shows that these morphisms preserve both
	$(\mathrm{C}^{\leq i-1}X)_\bullet$ and
	$(\mathrm{C}^{\leq i}X)_\bullet$ (this choice of bounds is optimal), and
	therefore induce
	% segment-id: o014.aljabr2.chapter8.dold-kan-b.q009
	\[ h_n:D_n\to D_{n+1}. \]

	Extend the definition of $s_k$ by zero to every $k\in\Z$, and thereby
	extend the definition of $h_n$ by zero to the case $n\leq i$. Notice that
	$n\leq i\implies D_n=0$, so this is the only reasonable choice. We claim
	that
	% segment-id: o014.aljabr2.chapter8.dold-kan-b.q010
	\[ \partial^D_{n+1}h_n+h_{n-1}\partial^D_n=\identity_{D_n}, \quad n\in\Z. \]
	This will show that $D$ is acyclic. First, applying
	\eqref{eqn:Dold-Kan-N-C-aux} gives
	% segment-id: o014.aljabr2.chapter8.dold-kan-b.q011
	\[ \partial^D_n = \sum_{j=0}^{n-i-2} (-1)^j d_j \;\bmod\; \mathrm{C}^{\leq i-1}(X)_{n-1}. \]
	Together with \eqref{eqn:simplicial-identity}, this gives, for
	$n\geq i+1$,
	% segment-id: o014.aljabr2.chapter8.dold-kan-b.q012
	\begin{align*}
		\partial^D_{n+1} h_n & = (-1)^{n-i-1} \sum_{j=0}^{n-i-1} (-1)^j d_j s_{n-i-1} \; \bmod\; \mathrm{C}^{\leq i-1}(X)_n \\
		& = (-1)^{n-i-1} \sum_{j=0}^{n-i-2} (-1)^j s_{n-i-2} d_j + \identity_{\mathrm{C}^{\leq i}(X)_n} \; \bmod\; \mathrm{C}^{\leq i-1}(X)_n \\
		& = -h_{n-1} \partial^D_n + \identity_{D_n}.
	\end{align*}
	The identity also holds trivially for $n\leq i$. This proves the claim.
\end{proof}

% segment-id: o014.aljabr2.chapter8.dold-kan-b.r001
\begin{remark}
	Naturally, one may ask for corresponding versions of the various theorems
	in this section, particularly the Dold--Kan correspondence, for the
	category of complexes $\cate{C}(\mathcal{A})$ or its subcategories
	$\cate{C}_{\leq0}(\mathcal{A})$ and
	$\cate{C}_{\geq0}(\mathcal{A})$. There are at least two simple
	approaches: reflection, and reversal or duality.
	\begin{enumerate}
		% segment-id: o014.aljabr2.chapter8.dold-kan-b.i006
		\item As explained in Remark~\ref{rem:cochain-vs-chain},
		$\cate{C}(\mathcal{A})$ and $\cate{Ch}(\mathcal{A})$ correspond by
		$X^n=X_{-n}$ and $d^n=\partial_{-n}$. All the results in this section
		can be transported to complexes through this correspondence. For
		example, the reflected version of Theorem~\ref{prop:Dold-Kan} is an
		adjoint equivalence between $\cate{C}_{\leq0}(\mathcal{A})$ and
		$\cate{s}\mathcal{A}$.
		% segment-id: o014.aljabr2.chapter8.dold-kan-b.i007
		\item Another approach is to consider the category of cosimplicial
		objects $\cate{cs}\mathcal{A}$. There are equivalences of categories
		% segment-id: o014.aljabr2.chapter8.dold-kan-b.q013
		\[ \cate{cs}\mathcal{A} \simeq \cate{s}\left(\mathcal{A}^{\opp}\right)^{\opp} \simeq \cate{Ch}_{\geq 0}(\mathcal{A}^{\opp})^{\opp} \simeq \cate{C}_{\geq 0}(\mathcal{A}), \]
		where the final step comes from reversing the arrows, as in
		Remark~\ref{rem:cochain-vs-chain}.
	\end{enumerate}
	Operations such as taking the normalized complex can likewise be
	translated into the category of complexes.
\end{remark}
